% =============================================================================
%  INTRODUCTION  &  NEUROCIRCUIT MODEL
%  Raw LaTeX body sections — no preamble, no \begin{document}
%  Custom macros assumed: \sectitle{}, \gene{}, \cit{}
% =============================================================================


% ─────────────────────────────────────────────────────────────────────────────
\sectitle{Introduction}
% ─────────────────────────────────────────────────────────────────────────────

Clinicians diagnose emotional disorders in categories---major depressive
disorder, borderline personality disorder, post-traumatic stress disorder,
attention-deficit/hyperactivity disorder---yet the patients who carry these
labels share a single, recurring failure: they cannot modulate the intensity,
duration, or expression of their emotional responses to match the demands of
the moment.  This mismatch between categorical nosology and dimensional
reality has plagued psychiatry for decades.  Comorbidity rates exceed 50\%
across mood, anxiety, and personality disorder diagnoses, and factor-analytic
studies consistently extract a general ``internalising'' dimension that
straddles every boundary the \emph{Diagnostic and Statistical Manual}
draws.\cite{faustino2020}  If the categories carve nature at its joints, the joints
should not bleed this freely.  The bleeding suggests that the categories are
wrong---not in their clinical utility, which remains indispensable for
treatment guidelines and insurance codes, but in their claim to reflect
discrete biological entities.  What biology offers instead is a continuum of
regulatory capacity that varies across individuals, contexts, and
developmental epochs, and that breaks down in characteristically different
ways depending on which circuits and genes are compromised.

The construct that cuts across every diagnostic joint is \emph{emotional
dysregulation}: a transdiagnostic deficit in the capacity to monitor,
evaluate, and modify emotional reactions in the service of adaptive goals.
Before Gross's foundational work, ``emotion regulation'' was a loosely
deployed umbrella term that conflated coping, defence mechanisms, mood
repair, and self-control without specifying which psychological operations
these constructs shared or how they differed.  Gross resolved this
ambiguity by proposing the process model of emotion regulation, which
decomposed the regulatory act into five sequential stages: situation
selection, situation modification, attentional deployment, cognitive
change, and response modulation.\cite{gross2002}  Each stage corresponds to a
distinct point in the emotion-generative trajectory at which a person can
intervene.  Situation selection and situation modification operate on the
external world before an emotional episode begins---choosing whether to
attend a funeral or rewriting the seating chart to avoid a hostile
relative.  Attentional deployment and cognitive change operate on internal
representations once the episode is underway---directing gaze away from
distressing stimuli or reinterpreting a critic's feedback as an
opportunity for growth.  Response modulation operates on the physiological
and behavioural output after the emotion has been fully generated---deep
breathing to slow heart rate, or suppressing a facial expression of anger.
This temporal ordering matters because antecedent-focused strategies, which
act before the full emotional response has been assembled, prove less
metabolically expensive and more effective than response-focused
strategies, which attempt to suppress an emotion already in flight.\cite{gross2002}
The process model thus transformed emotion regulation from a
folk-psychological concept into a mechanistically specified sequence with
measurable decision points, each amenable to experimental manipulation and
each mapping onto identifiable neural substrates.

John and Gross sharpened this five-stage architecture into two prototypical
strategies that have since anchored thousands of empirical
studies.\cite{john2004}  Cognitive reappraisal---reinterpreting a situation to alter
its emotional meaning---reduces negative affect, preserves social
connection, and costs little in working memory.  Expressive
suppression---inhibiting the outward display of an emotion without changing
the internal experience---fails to reduce subjective distress, impairs
episodic memory for the triggering event, and erodes interpersonal rapport
over repeated interactions.  The asymmetry between these two strategies is
not merely quantitative; it is qualitative.  Reappraisal alters the
emotional trajectory early, before autonomic and endocrine cascades reach
full amplitude, while suppression leaves the full physiological response
intact and adds a second metabolic burden---the effortful inhibition of
expressive behaviour---on top of it.  Habitual reappraisers show lower
sympathetic nervous system activation during social stress, lower cortisol
reactivity during laboratory challenges, and faster cardiovascular recovery
after negative emotional episodes.  Habitual suppressors show the reverse
on every measure.\cite{john2004}  Individual differences in habitual strategy use
therefore predict well-being more reliably than the frequency or intensity
of negative life events: people who habitually reappraise report higher
life satisfaction, fewer depressive symptoms, and stronger social bonds,
while people who habitually suppress report lower satisfaction, more
symptoms, and progressive social withdrawal.  The reappraisal--suppression
axis offers a dimensional metric of regulatory health that transcends
diagnostic categories.  A patient diagnosed with PTSD who habitually
reappraises may function better than a patient with no formal diagnosis who
habitually suppresses---a paradox that categorical nosology cannot
accommodate but that dimensional models predict and explain.

Building on this dimensional logic, Linehan's biosocial theory located the
developmental origins of severe dysregulation at the intersection of
biological vulnerability and environmental invalidation.\cite{crowell2009}  Children
born with heightened emotional sensitivity, rapid emotional escalation, and
slow return to baseline---a temperamental triad mediated by amygdala
reactivity and prefrontal maturation rates---develop adaptive regulation
only when their caregiving environment acknowledges, labels, and coaches
emotional experience.  When the environment instead dismisses, punishes, or
ignores emotional expression, the child never acquires the cognitive
reappraisal strategies that Gross's model identifies as adaptive.  The
result is a regulatory repertoire dominated by suppression, avoidance, and
impulsive action---strategies that provide momentary relief at the cost of
long-term escalation.  Crowell and colleagues formalised this account as a
developmental cascade: heritable neurocognitive vulnerabilities interact
with invalidating environments across sensitive periods to produce the
affective instability, identity disturbance, and interpersonal chaos
characteristic of borderline personality disorder.\cite{crowell2009}  Crucially,
however, the same cascade---with different weightings of vulnerability and
invalidation---generates the dysregulation observed in PTSD, recurrent
depression, ADHD, and substance-use disorders.  In PTSD, the vulnerability
is a sensitised threat circuit and the invalidation is the trauma itself;
in ADHD, the vulnerability is a sluggish reward circuit and the
invalidation is an environment that demands sustained effortful control; in
substance-use disorders, the vulnerability is an over-responsive
dopaminergic system and the invalidation is early exposure to substances
during a period of prefrontal immaturity.  The biosocial model is therefore
not a theory of one diagnosis; it is a theory of one mechanism expressed
across many diagnoses.

This transdiagnostic convergence motivated the National Institute of Mental
Health to abandon categorical funding priorities in favour of the Research
Domain Criteria (RDoC) framework.\cite{insel2010}  RDoC organises psychopathology
along five dimensional constructs---negative valence systems, positive
valence systems, cognitive systems, social processes, and
arousal/regulatory systems---each analysed across seven units of analysis
that span genes, molecules, cells, circuits, physiology, behaviour, and
self-report.  Emotional dysregulation sits at the intersection of the
negative valence and arousal/regulatory domains, precisely where
circuit-level and genetic evidence converge most densely.  Within the
negative valence domain, dysregulation maps onto the ``sustained threat''
and ``frustrative nonreward'' constructs; within the arousal/regulatory
domain, it maps onto the ``arousal'' and ``circadian rhythms'' constructs
that govern the temporal dynamics of emotional responding.  By replacing
the question ``Which disorder does this patient have?'' with ``Which
regulatory dimensions are disrupted, at which units of analysis, and to
what degree?'', RDoC reframes emotional dysregulation as a measurable,
mechanistically decomposable target rather than a loose clinical
descriptor.  The framework also encourages researchers to study the full
range of regulatory variation---from resilient individuals who regulate
effortlessly through subclinical populations who struggle under stress to
clinical populations who cannot regulate at all---rather than comparing
categorical ``cases'' with categorical ``controls'' and discarding the
dimensional information that lies between them.

Phenotype ontologies offer a complementary strategy for escaping categorical
constraints.  The Human Phenotype Ontology (HPO) term
\texttt{HP:0100851}---Abnormal emotional state---subsumes twelve descendant
phenotypes including anhedonia (\texttt{HP:0012154}), irritability
(\texttt{HP:0000737}), dysphoria (\texttt{HP:0100020}), emotional lability
(\texttt{HP:0000712}), anger (\texttt{HP:0031473}), and apathy
(\texttt{HP:0000741}).\cite{kohler2017}  Each descendant maps to distinct neural
circuit signatures, genetic architectures, and pharmacological response
profiles, yet all share the parent node.  This hierarchical structure
mirrors the empirical reality that dysregulation is not one phenotype but a
family of phenotypes organised by severity, valence, and temporal dynamics.
Anhedonia involves a failure to generate positive affect; irritability
involves a failure to inhibit negative affect; emotional lability involves
rapid oscillation between affective states; apathy involves a failure to
generate any affect at all.  These facets can co-occur, dissociate, and
respond differently to treatment within the same individual.  Mapping
clinical observations onto HPO terms enables cross-study harmonisation,
large-scale phenome-wide association studies, and the precise specification
of which \emph{facet} of dysregulation a given gene, circuit, or
intervention targets.

The present review exploits this dimensional, multi-level architecture.  We
begin by mapping the neural circuitry that implements---and fails to
implement---emotion regulation, drawing on converging evidence from
structural and functional neuroimaging, lesion studies, and causal
perturbation experiments.  We then trace the genetic architecture that
tunes this circuitry, moving from genome-wide association signals through
gene-expression landscapes to rare-variant convergence on shared synaptic
pathways.  Next, we examine developmental trajectories in which genetic risk
and circuit maturation interact with environmental exposures to crystallise
stable patterns of dysregulation.  Finally, we evaluate therapeutic
interventions---pharmacological, psychotherapeutic, and
neuromodulatory---through the lens of the circuits and genes they engage.
Throughout, we anchor our analysis in the RDoC framework and HPO phenotype
hierarchy, treating emotional dysregulation not as a symptom of a category
but as a measurable dimension that demands mechanistic explanation at every
level of biological organisation.


% ─────────────────────────────────────────────────────────────────────────────
\sectitle{Neurocircuit Model of Emotional Dysregulation}
% ─────────────────────────────────────────────────────────────────────────────

The prefrontal--amygdala axis has dominated neurocircuit models of emotion
regulation for two decades, yet the field's understanding of \emph{how}
this axis fails has undergone a fundamental revision.  Early formulations
cast the amygdala as a threat-detection alarm and the prefrontal cortex as
its top-down brake: dysregulation, on this account, arose from an
overactive amygdala that overwhelmed prefrontal inhibition.  This ``see-saw''
model was intuitive, clinically appealing, and wrong in its central claim.
Silverman and colleagues overturned the simple story in a large-sample
(\emph{N}\,>\,500) study that parsed neuroticism---the personality trait
most strongly and consistently associated with emotional
dysregulation---into its constituent circuit signatures.\cite{silverman2019}  They found
that neuroticism predicted neither amygdala hyperactivation nor prefrontal
hypoactivation in isolation; instead, it tracked \emph{deficient functional
connectivity} between the amygdala and ventromedial prefrontal cortex
(vmPFC).  The amygdala generated threat signals at normal amplitude, but
those signals failed to propagate to the prefrontal regions that would
contextualise, reappraise, and ultimately down-regulate them.
Dysregulation, this finding revealed, is a disconnection disorder, not an
excitation disorder---a failure of communication between nodes, not a
failure of any single node.

Kredlow and colleagues extended this connectivity-centric view to PTSD by
decomposing prefrontal--amygdala interactions into three temporally distinct
phases: fear acquisition, fear extinction, and fear regulation.\cite{kredlow2021}
During acquisition, patients with PTSD showed intact amygdala
signalling---they learned the threat contingency as rapidly as
trauma-exposed controls, and their amygdala BOLD responses to conditioned
stimuli were indistinguishable from those of healthy participants.  During
extinction, however, vmPFC failed to suppress the acquired fear response:
the safety signal was encoded but could not override the threat
representation.  During active regulation---when participants were
instructed to reappraise the threatening stimulus---dlPFC failed to recruit
vmPFC to implement the cognitive transformation that reappraisal demands.
This three-phase dissection clarified a clinical puzzle that had resisted
explanation for years: why patients with PTSD \emph{can} learn new safety
signals in the controlled laboratory environment yet \emph{cannot} deploy
them in ecologically valid contexts where competing demands strain
prefrontal resources.  The acquisition circuit works; the extinction and
regulation circuits do not.  The therapeutic implication follows
directly---interventions must target connectivity during extinction and
regulation, not amygdala reactivity during acquisition, and
exposure-based therapies must be augmented with strategies that strengthen
the prefrontal--amygdala pathway rather than simply re-presenting the
feared stimulus.

Depression presents a complementary pattern that further fragments the
monolithic ``PFC dysfunction'' narrative.  Pizzagalli and Roberts reviewed
convergent evidence from task-based fMRI, resting-state connectivity
analyses, and computational reinforcement-learning models to show that
prefrontal disruption in depression is subtype-specific.\cite{pizzagalli2021}  Anhedonic
depression---characterised by diminished pleasure, reduced motivation, and
blunted reward learning---maps to reduced connectivity between medial PFC
and ventral striatum, reflecting an impaired reward-prediction error signal
that starves the cortex of the positive-valence information needed to guide
goal-directed behaviour.  Anxious depression---characterised by
hypervigilance, somatic tension, and ruminative worry---maps to exaggerated
connectivity between the dorsal anterior cingulate cortex (dACC) and the
amygdala, reflecting an overtuned threat-monitoring loop that biases
attention toward negative stimuli at the expense of neutral and positive
alternatives.  Melancholic depression---characterised by psychomotor
retardation, diurnal mood variation, and vegetative signs---maps to
reduced dlPFC activation during cognitive control tasks, reflecting a
depleted regulatory resource that can neither reappraise threat nor amplify
reward.  No single prefrontal abnormality captures ``depression''; the
diagnostic label masks at least three circuit-level disorders, each with
distinct pharmacological and psychotherapeutic response profiles.  This
heterogeneity explains decades of inconsistent neuroimaging findings and
argues forcefully for the dimensional, circuit-specified approach that RDoC
demands.  It also suggests that treatment selection should be guided not by
the syndromal diagnosis but by the circuit subtype---a strategy that
emerging imaging-guided treatment algorithms are beginning to test.

Beyond mood and anxiety disorders, the prefrontal--subcortical axis also
governs the dysregulated reward pursuit that characterises addictive
behaviours.  Brand and colleagues proposed the Interaction of
Person--Affect--Cognition--Execution (I-PACE) model, which has since
attracted over 1,600 citations, to formalise how fronto-striatal imbalance
drives compulsive engagement with substances and behavioural rewards
alike.\cite{brand2019}  In their framework, affective and cognitive responses to
triggers interact within a prefrontal--striatal loop: ventral striatal
signals encode the incentive salience of a cue, while dlPFC and inferior
frontal gyrus (IFG) compute the long-term costs of acting on that cue.
When prefrontal inputs weaken---through chronic stress, substance-induced
neurotoxicity, or trait-level impulsivity that reflects developmental
prefrontal immaturity---the striatal signal wins, and the person pursues
the reward despite foreseeing its negative consequences.  The I-PACE model
thus extends the disconnection framework from threat regulation
(amygdala--PFC) to reward regulation (striatum--PFC), revealing a shared
architectural principle: dysregulation arises whenever a subcortical signal
is generated at normal or elevated intensity but the prefrontal circuit
that should contextualise and constrain that signal is structurally or
functionally disconnected.  This principle unifies threat-related and
reward-related forms of dysregulation under a single computational
logic---top-down contextualisation fails, and the subcortical signal drives
behaviour unopposed.

Recent causal evidence has elevated this correlational architecture to
mechanistic status.  Niu and colleagues used concurrent transcranial
magnetic stimulation and functional MRI (TMS--fMRI) to demonstrate that
single-pulse stimulation of the left dlPFC produced immediate, measurable
changes in amygdala blood-oxygen-level-dependent (BOLD) signal in healthy
participants.\cite{niu2026}  The same stimulation protocol, applied to
individuals with subthreshold depression, \emph{failed} to modulate
amygdala activity: the dlPFC was activated, but the downstream signal did
not arrive.  This causal probe confirmed what correlational connectivity
analyses had inferred: the dlPFC-to-amygdala regulatory pathway is not
merely associated with mood state but is causally necessary for effective
emotional down-regulation, and this causal link degrades before clinical
thresholds are crossed.  Subthreshold depression, on this evidence, is
already a circuit disorder, not merely a risk factor for one.  The clinical
implication is that intervention at the subthreshold stage---when the
pathway is weakened but not yet severed---may prevent the transition to
full syndromal disorder, a possibility that current neurostimulation trials
are actively exploring.

Meiering and colleagues approached the same circuitry from the
trait-vulnerability direction, using high-resolution fMRI during an
emotional face-processing task to show that neuroticism predicts increased
\emph{functional} connectivity among the amygdala, hippocampus, and medial
PFC during the processing of negative emotional stimuli.\cite{meiering2026}  At first
glance, this finding appears to contradict the Silverman result of
\emph{deficient} amygdala--vmPFC connectivity in high-neuroticism
individuals.  The resolution lies in the critical distinction between
regulatory connectivity and reactive connectivity.  Silverman measured
connectivity during explicit emotion regulation---participants were
instructed to down-regulate their emotional response; Meiering measured
connectivity during passive emotional processing---participants simply
viewed the stimuli without regulatory instruction.  Neuroticism, therefore,
involves a double dissociation: heightened amygdala--hippocampal--PFC
coupling during the initial encoding of emotional events, which amplifies
the subjective salience of negative experience and strengthens the
contextual memory trace that binds the negative emotion to its situational
triggers, and weakened amygdala--vmPFC coupling during the subsequent
attempt to regulate that experience, which prevents effective
down-regulation.  The neurotic brain, in short, amplifies the emotional
signal through enhanced reactive connectivity and then disconnects the
regulatory brake through impaired top-down control.  This double
dissociation explains why highly neurotic individuals report both more
intense negative emotions \emph{and} more difficulty managing those
emotions---features that a simple hyperactivation or hypoactivation model
cannot capture.

These macro-circuit findings acquire finer anatomical grain when mapped
onto the Allen Human Brain Atlas, which parcellates the human amygdala into
88 sub-regions spanning the lateral nucleus, basal nucleus, accessory basal
nucleus, central nucleus, cortical nuclei, and their laminar
subdivisions.\cite{hawrylycz2012}  The lateral nucleus receives the densest cortical
input---from auditory, visual, and polymodal association cortices---and
serves as the primary gateway for sensory information entering the
amygdala.  The central nucleus projects to brainstem effector regions,
including the periaqueductal grey, parabrachial nucleus, and dorsal vagal
complex, that generate the autonomic, endocrine, and motor components of
the emotional response.  The basal nucleus bridges these input and output
streams while maintaining dense reciprocal connections with vmPFC,
orbitofrontal cortex, and the nucleus accumbens, positioning it as the
critical hub through which prefrontal cortex exerts its regulatory
influence.  The connectivity deficits reported by Silverman, Kredlow, and
Niu likely localise to the basal nucleus--vmPFC pathway; the reactive
connectivity amplification reported by Meiering likely engages the lateral
nucleus--hippocampal pathway.  Future studies that combine ultra-high-field
imaging (7T and above) with atlas-informed parcellation will be positioned
to test this sub-regional hypothesis directly, moving the field from
treating the ``amygdala'' as a monolithic computational node to recognising
it as a differentiated structure whose sub-circuits serve distinct
regulatory, reactive, and output-generation functions.

How, then, do these circuit disruptions relate to the adaptive regulation
that Gross's process model describes?  Troy and colleagues addressed this
question by proposing an affect-regulation framework for psychological
resilience.\cite{troy2022}  In their model, positive reappraisal---the cognitive
strategy of reinterpreting a negative event as an opportunity for growth,
learning, or meaning---serves as the central mechanism linking prefrontal
circuit integrity to adaptive outcomes.  Resilient individuals, Troy and
colleagues showed, do not experience less adversity and do not generate
weaker amygdala responses to that adversity; rather, they recruit vmPFC and
dlPFC more effectively during reappraisal, which transforms the valence of
the amygdala signal before it propagates to downstream effector circuits.
The reappraised signal still carries information about the threatening
stimulus, but that information is now contextualised by a narrative of
agency, growth, or meaning that reduces its autonomic impact and frees
attentional resources for goal-directed action.  Resilience, on this
account, is not a fixed trait that resides in a brain region; it is a
dynamic process that emerges from the efficiency of a circuit---precisely
the amygdala--vmPFC--dlPFC circuit that Silverman, Kredlow, and Niu have
shown to be compromised in neuroticism, PTSD, and subthreshold depression.
The Troy framework thereby closes the loop between clinical deficit and
adaptive function: the same circuit that, when intact, enables reappraisal
and resilience, when degraded, produces the emotional dysregulation observed
across diagnostic categories.

Together, these findings converge on a unified neurocircuit model of
emotional dysregulation that has three core properties.  First,
dysregulation is a connectivity disorder: the critical failure lies not in
the amplitude of subcortical threat or reward signals but in the strength,
timing, and fidelity of prefrontal modulation of those signals.  Second,
the model is subtype-specific: amygdala--vmPFC disconnection drives
threat-related dysregulation (neuroticism, PTSD, generalised anxiety),
mPFC--ventral striatal disconnection drives reward-related dysregulation
(addiction, anhedonic depression), and dACC--amygdala hyperconnectivity
drives the anxious hypervigilance that cuts across anxiety and depressive
disorders.  Third, the model is dimensional: circuit degradation manifests
on a continuum from subclinical trait vulnerability through prodromal states
to full syndromal disorder, with no categorical boundary demarcating health
from illness.  Each position on this continuum corresponds to a measurable
degree of connectivity strength that predicts regulatory capacity,
treatment response, and longitudinal trajectory.

This circuit architecture raises an immediate mechanistic question.  If
dysregulation is fundamentally a connectivity disorder, what determines the
strength, myelination, and receptor density of these long-range
prefrontal--subcortical connections?  Synaptic density, white-matter
integrity, neurotransmitter receptor distribution, and monoamine
metabolism all shape the efficiency of cortical--subcortical
communication, and twin studies estimate that 40--60\% of the variance
in prefrontal--amygdala connectivity is heritable.\cite{polderman2015}  The
prefrontal--amygdala pathway is among the last to myelinate during
development---full maturation extends into the mid-twenties---and among
the most sensitive to experience-dependent synaptic pruning during
adolescence.  This protracted maturational window creates ample
opportunity for genetic variants to shift developmental trajectories
toward or away from dysregulation, and for gene--environment interactions
to amplify or buffer those genetic effects.  Understanding which genes
tune this circuitry, how their expression changes across developmental
epochs, how they interact with environmental exposures, and where they
converge in biological pathways is therefore essential for completing the
mechanistic account that the neurocircuit model demands.  We turn to
this genetic architecture next.
